%\vspace{-2mm}
\section{Experiments on COLIEE dataset} \label{sec:coliee}
%\vspace{-3mm}

\begin{table}[t]
    \centering
    %\small
    \small
    \setlength{\tabcolsep}{5pt}
    \begin{tabular}{p{0.1\textwidth}cccc}
        \toprule
        \textbf{Model} &  \multicolumn{2}{c}{\textbf{Statutes}}&  \multicolumn{2}{c}{\textbf{Precedents}}\\
         &  \ilpcsr &  \coliee &  \ilpcsr &  \coliee \\ \midrule
         BM25 (5-gram)&  16.98&  54.49&  33.29& 30.61\\
         Event-GNN&  28.67& - &  12.08& 11.48 \\
         Para-GNN&  20.72&  17.64&  24.54& 21.24\\
 Para-GNN (summaries)& & -& &18.52\\ \midrule
         Para-GNN + BM25&  28.10 & 55.96&  \textbf{36.93} & \textbf{34.52}\\
Event-GNN + BM25& 33.87& -& 34.45&34.51\\ 
 Para-GNN (summaries) + BM25& \textbf{36.17} & -& 36.35&30.25\\ \bottomrule
    \end{tabular}
    \vspace{-3mm}
    \caption{Results of the best methods on \coliee{} datasets compared to ILPCSR. All results are in terms of macro-F1@$K$. Event-GNN and summaries could not be run for \coliee{} statutes since the queries are too small for meaningful events or summaries.}
    \label{tab:coliee}
    \vspace{-6mm}
\end{table}

All our observations are drawn over the \ilpcsr{} dataset constructed in this work. We would like to verify whether the above trends are also seen on other legal datasets/jurisdictions. Since no current dataset allows the identification of both statutes and precedents together from the same query, we have to work with two separate datasets for LSR and PCR. We choose to work on the well-known COLIEE datasets~\cite{li2024coliee2024}. 
%Specifically, we use the COLIEE 2024 Task 1 (Legal Precedent Retrieval) and Task 3 (Statute Law Retrieval) datasets for the tasks of LSR and PCR respectively. 

\noindent\underline{\textbf{LSR dataset:}} We use the \texttt{COLIEE 2024 Task 3 (Statute Law Retrieval)} dataset consisting of 746 statutes from \textit{Japanese law}, and 554 queries. Note that the queries here are typically one or two sentences long, asking specifically about the laws. 
In contrast, the queries in \ilpcsr{} are real-life cases, which makes the setting more practical and challenging. 
We opted for this dataset since other existing datasets in English (ECHR2021 and ILSI) have too less statutes  (66 and 100 respectively) to be evaluated in the retrieval setup. 

\noindent \underline{\textbf{PCR dataset:}} We use the \texttt{COLIEE 2024 Task 1 (Legal Precedent Retrieval)} dataset consisting of 1,678 queries and 5,529 precedent candidates, all of which are real-life case judgments from \textit{Canadian Federal law}. This setting is similar to the queries in \ilpcsr{}. 

\noindent \underline{\textbf{Results:}}
We choose some of the methods that performed highly over the \ilpcsr{} dataset, and applied these methods on the COLIEE datasets.
The results on \coliee{} vis-a-vis \ilpcsr{} are presented in Table~\ref{tab:coliee}.
The trends on the \coliee{} dataset are almost similar to what we observed for \ilpcsr{}. 
The only difference being that, for \coliee{}, even for LSR, lexical methods such as BM25 perform the best (whereas semantic methods outperformed lexical approaches for LSR in \ilpcsr{}). 
This difference is possibly because the queries of \coliee{} are short sentences, asking directly about the statutes, whereas for \ilpcsr{} the queries are real-world long case judgments. Both for LSR and PCR, we see improvements when using an ensemble setup for \coliee{} as well. The improvement is limited in case of statutes, possibly because the performance of Para-GNN is poor. This is possibly because short queries do not have enough structure for the GNN to exploit. For precedents, where the semantic methods perform better, the improvement obtained by ensembling is high. 
This agrees with the trend we see on \ilpcsr{} (see Table~\ref{tab:coliee}). Finally, for both \ilpcsr{} and \coliee{}, in the case of PCR, we have observed that using summaries does not perform as well as using the full texts. We observe the same key findings as \ilpcsr. 
